problem:  
Let \((P,\mathcal{H},g)\) be a complete, connected, equiregular sub‑Riemannian manifold endowed with a smooth, free, proper, and isometric action of a compact Lie group \(G\) preserving the Popp measure \(\mu_P\), the horizontal distribution \(\mathcal{H}\), and its metric \(g\). Assume the horizontal connection \(\nabla^{\mathcal{H}}\) allows a \(G\)-invariant decomposition of \(TP\) into horizontal and vertical bundles and defines an equivariant connection 1‑form \(\eta\) with curvature \(\Omega=d\eta+ \frac12[\eta,\eta]\).  
Let \(\pi:P\to M=P/G\) be the quotient map, and denote by \((\mathcal{H}_{\mathrm{red}},g_{\mathrm{red}})\) the reduced sub‑Riemannian structure induced on \(M\).

Every normal sub‑Riemannian geodesic on \(P\) is an integral curve of the Hamiltonian flow of \(H_P(p)=\tfrac12 g^{-1}(p|_{\mathcal{H}},p|_{\mathcal{H}})\) on \(T^*P\).  Let \(\mathbf{J}:T^*P\to \mathfrak{g}^*\) be the cotangent‑lifted moment map of the \(G\)‑action, and consider a regular value \(\mu\in\mathfrak{g}^*\).  The reduced Hamiltonian \(H_{P,\mu}\) on \(\mathbf{J}^{-1}(\mu)/G_\mu\) projects to the Hamiltonian \(H_M\) of the reduced system on \(T^*M\).  

For a normal geodesic \(\gamma\) of \(H_M\), denote by \(\Lambda_\gamma(t)\subset T_{\lambda(t)}T^*M\) its Jacobi curve and by \(\tilde\Lambda_{\tilde\gamma}(t)\subset T_{\tilde\lambda(t)}T^*P\) the lifted Jacobi curve along its horizontal lift \(\tilde\gamma\) that satisfies \(\mathbf{J}(\tilde\lambda(t))=\mu\).

**Refined Open Problem (Jacobi curve under equivariant symplectic reduction):**  
Assume \(G\) acts freely and \(\mu\) is a regular value of \(\mathbf{J}\).  
Is it true that
\[
\Lambda_\gamma(t) \;\cong_{\mathrm{Sp}}\; 
\bigl(\tilde\Lambda_{\tilde\gamma}(t) \cap T_{\tilde\lambda(t)}(\mathbf{J}^{-1}(\mu))\bigr)\big/\!\!\sim,
\]
that is, that the Jacobi curve of the reduced geodesic is (up to symplectic isomorphism) the **symplectic quotient** of the ambient Jacobi curve restricted to the moment‑map level set?  

Further, can the curvature invariants of \(\Lambda_\gamma\) (such as the eigenvalues of the Agrachev–Zelenko curvature operator) be expressed in terms of those of \(\tilde\Lambda_{\tilde\gamma}\) together with the curvature tensor \(\Omega\) of the principal \(G\)-connection?

In other words, does there exist a canonical “Jacobi reduction formula”
\[
\mathcal{R}^M = \pi_*\!\big(\mathcal{R}^P|_{\mathbf{J}^{-1}(\mu)}\big) + \Psi(\Omega,\mu),
\]
for some tensorial correction term \(\Psi(\Omega,\mu)\) determined purely by the connection curvature and the moment level \(\mu\)?  

This conjectural formula would generalize the O’Neill tensor in Riemannian submersion theory to the non-holonomic, symplectically reduced setting.

Why is it a "good" problem:  
This sharper formulation poses a precise new question: whether **Jacobi curves—central symplectic invariants of sub‑Riemannian geometry—descend canonically through Marsden–Weinstein reduction** and how their canonical curvature operators transform. The problem lies at a deep intersection of **symplectic geometry** (moment–map reduction) and **sub‑Riemannian curvature theory** (Agrachev–Zelenko formalism). Its statement is simple—“does the Jacobi curve reduce symplectically?”—yet conceptually profound, as its resolution would unify the geometric structure of Hamiltonian reduction with non‑holonomic curvature. Establishing such a formula would create a new bridge between curvature invariants, group symmetries, and symplectic reduction, analogous to (but deeper than) classical Riemannian submersion formulas, thus fully qualifying as a “good” research problem.